%! Unit 1 — Functions and Change — Unit Cover Page
\documentclass[10pt]{article}
\usepackage{precalculus-article}
\usepackage{precalculus-boxes}
\usepackage{ltablex}
\keepXColumns

\begin{document}
\thispagestyle{empty}

% Full-bleed plum banner
\begin{tikzpicture}[remember picture, overlay]
  \fill[plum] (current page.north west)
    rectangle ([yshift=-1.15in]current page.north east);
\end{tikzpicture}
\vspace{-0.82in}

\begin{center}
  {\color{white}\textbf{\LARGE Precalculus}}\\[6pt]
  {\color{lilacmid}\Large\textbf{Unit 1: Functions and Change}}\\[4pt]
  {\color{white}\small 8 Lessons + Unit Test \quad|\quad Class Periods: 18--22}
\end{center}

\vspace{0.24in}

\begin{tcolorbox}[colback=goldbg, colframe=goldacc, boxrule=0.8pt, arc=2mm,
    left=3mm, right=3mm, top=2mm, bottom=2mm]
\small\textbf{Unit Essential Questions:}
\begin{itemize}[leftmargin=1.4em, itemsep=2pt, topsep=2pt]
  \item How do we describe the way one quantity changes with respect to another?
  \item What do rates of change, transformations, composition, and inverses reveal
        about how a function behaves?
\end{itemize}
\end{tcolorbox}

\vspace{0.12in}

\begin{tocbox}
\small
\renewcommand{\arraystretch}{1.42}
\begin{tabularx}{\linewidth}{@{} c >{\bfseries}X l @{}}
\rowcolor{lilac}
\textbf{\#} & \textbf{Lesson Title} & \textbf{Content Source} \\
\midrule
1.0 & Introduction to Functions and Change & New \\
1.1 & Function Fundamentals & New \\
1.2 & Change in Tandem & CED 1.1 \\
1.3 & Rates of Change & CED 1.2 + 1.3 \\
1.4 & Transformations of Functions & CED 1.12 \\
1.5 & Composition of Functions & CED 2.7 \\
1.6 & Inverse Functions & CED 2.8 \\
1.7 & Function Model Selection and Construction & CED 1.13 + 1.14 \\
\midrule
\rowcolor{lilac}
    & \textbf{Sample Test} & \\
\end{tabularx}
\end{tocbox}

\vspace{0.12in}

\begin{skillbox}[Mathematical Practices --- Unit 1 Focus]{goldbox}
\renewcommand{\arraystretch}{1.35}
\begin{tabularx}{\linewidth}{@{} >{\leavevmode\bfseries\color{plum}}p{0.12\linewidth} X @{}}
MP 1.C & Construct new functions using transformations, compositions, inverses, or regressions. \\
MP 2.A & Identify information from graphical, numerical, analytical, and verbal representations. \\
MP 2.B & Construct equivalent graphical, numerical, analytical, and verbal representations. \\
MP 3.A & Describe the characteristics of a function with varying levels of precision. \\
MP 3.C & Support conclusions or choices with a logical rationale or appropriate data. \\
\end{tabularx}
\end{skillbox}

\vspace{0.12in}

\begin{remindbox}
\small
One idea carries the whole unit: a function is a rule that pairs each input with
exactly one output, and everything else here is a way of \textbf{describing} or
\textbf{building} such a rule. Lessons 1.0--1.3 describe how a function changes
--- increasing, decreasing, concavity, average rate of change. Lessons 1.4--1.6
build new functions from old ones --- shifting and stretching, feeding one
function into another, and reversing one. Lesson 1.7 puts both halves to work:
choosing a model that fits the data and stating what it assumes.
\end{remindbox}

\end{document}
